\documentclass[12pt]{article}
\usepackage{lmodern}
\usepackage[T1]{fontenc}
\usepackage[letterpaper,margin=0.75in,bottom=0.75in]{geometry}
\usepackage{cancel}
\usepackage{graphicx}
\usepackage{braket}
\usepackage{latexsym,amssymb,amsmath}
\usepackage{pdfpages}
\usepackage{xcolor}
\usepackage{capt-of}
\usepackage{cite}


\begin{document}

%%%%%%%%%%%%%%%%%%%%%%%%%%%%%%%%%%%%%%%%%%%%%%%%%%%%%%%%%%%%%%%%
%%%%%%%%%%%%%%%%%%%%%%%%%%%%%%%%%%%%%%%%%%%%%%%%%%%%%%%%%%%%%%%%
\title{Proposal for Revising the AB Physics Major \\
Department of Physics and Astronomy}
%%%%%%%%%%%%%%%%%%%%%%%%%%%%%%%%%%%%%%%%%%%%%%%%%%%%%%%%%%%%%%%%
%%%%%%%%%%%%%%%%%%%%%%%%%%%%%%%%%%%%%%%%%%%%%%%%%%%%%%%%%%%%%%%%

\maketitle

\section*{Summary}

\noindent
In this section we present a summary of our proposal to revise the
undergraduate {\bf AB Physics} major.  Further details are included in
the sections that follow.  We have also
provided, along with this brief, a longer write-up that describes our
complete plan for revising the curriculum, including the BS in Physics
and BS in Applied Physics majors.\\[5pt]

\noindent
This proposal modifies the AB physics major in several complementary ways, mostly aimed at improving
student instruction, informed by our experience with the current program and observed shortcomings.
This proposal removes the requirement for a fifth quarter in the honors introductory physics
sequence (PHY 9HE), which will be discontinued.  This will allow students following this sequence to
start on other requirements earlier.  The proposal updates the math requirements to include
additional new options.  The newer math courses include all of the material needed by our majors,
and are already being accepted as equivalent on an ad-hoc basis.  This proposal simply makes this
policy explicit.  The requirement for PHY 102 (one unit) has been dropped, and PHY 40
(three units), already needed as a prerequisite of PHY 104A, is now
required explicitly.  Experience has shown that it is not feasible to
introduce computational physics with one unit.  PHY 40 is a completely
redesigned course, aimed to introduce computational physics, with no
physics or math prerequisites, and has been assigned an appropriate
number of units.  We have added PHY 117 and 118 and two new
astronomy courses (PHY 156 and PHY 158) to the electives, and students
now choose at least two electives.  We have also eliminated the
requirement for an additional upper-division physics course of any
kind.\\[5pt]

\noindent
The current AB Physics major requires a minimum of 85 units, which
exceeds the unit cap (80) by five units, and 35 upper division units,
one fewer than the 36 required by L\&S Regulation 34(A)(1).  The
proposed requirements can be completed with a minimum of 86 units, 36
of them upper division.  No path through the AB can meet both
limits, because the preparatory coursework alone requires 50 units.
The department therefore requests an exemption under L\&S Regulation
33(D), presented in the Required Units section below.

\newpage

\section*{Revision of Existing Majors and Minors}
{\it Please submit proposals for revision of an individual major or minor via email to Associate Director Mary Vasquez at \texttt{mpvasquez@ucdavis.edu}}.\\[5pt]

\noindent
Here we present our proposal to revise the undergraduate {\bf AB
  Physics} major.  For convenience, we include the specific prompts
from the guidance in ``Establishment, Revision, and Discontinuation of
Undergraduate Academic Degree Programs'', with matching section
numbers and prompt letters.

\section*{1. Approvals}

\noindent
{\bf Prompt a:} {\it Include all faculty votes and approval letters as required by PPM 200-25.}\\[5pt]

\noindent
This proposal has been developed over more than five years within the
Department of Physics and Astronomy, led by the undergraduate
curriculum committee.  The proposal was subjected to an extensive
vetting process that involved assigned department readers who were not
involved in formulating the initial plan.  {\bf The proposal was
unanimously supported by the department in a December 2021 vote.}

\section*{2. Background and Rationale}

\noindent
{\bf Prompt a:} {\it Explain all proposed revisions to the currently approved program.}\\[5pt]

\noindent
The updated AB physics major requirements include the following changes:
\begin{enumerate}
\item PHY 9HE is removed from the Honors Physics sequence, and the
  course will be discontinued.
\item The math requirements have been updated to include additional
  new options available.
\item The PHY 102 requirement has been dropped, and PHY 40, already needed as a prerequisite of PHY 104A, is now required explicitly.
\item PHY 117 and 118 and new Astrophysics courses (PHY 156 and 158) have been added as electives.
\item At least two electives are now required, instead of one.
\item The requirement for one additional upper-division physics course has been removed.
\item The introductory physics requirement has been restructured:  PHY
  9A or 9HA, followed by the rest of either the PHY 9 sequence
  (9B--9D) or the PHY 9H sequence (9HB--9HD).  A student may switch
  between the two sequences only once, immediately after the first
  course.
\item Minor clarification:  PHY 185 and 190 are recommended.
\end{enumerate}

\noindent
{\bf Prompt b:} {\it Explain the rationale for why the proposed revisions are necessary.}\\[5pt]

\noindent
Combined with already implemented schedule changes for other course
offerings, eliminating the PHY 9HE requirement allows students following
the honors physics sequence to start on other requirements starting in
the winter quarter of their sophomore year.  This allows them to
balance their physics requirements more evenly throughout their four
years.  An additional benefit of the already implemented schedule changes,
including twice per year offerings of PHY 104A (Introduction to
Mathematical Methods in Physics), is that the course load for
transfer students is more balanced as well.\\[5pt]

\noindent
The PHY 102 (Computational Laboratory in Physics) was generally taught
as a 3-4 unit course despite being only one-unit.  Experience has
shown that a one-unit course is insufficient to introduce
computational physics.  PHY 102 will be discontinued.  The new
requirements add PHY 40 (Introduction to Computational Physics).
Computational Physics continues to grow in importance in our field,
and PHY 40 is a prerequisite of PHY 104A, a required course.\\[5pt]

\noindent
We have been approving students' requests to substitute newly available
math courses (MAT 27A, 67, 27B) where appropriate for our current
requirements (MAT 22A and 22B).  This proposal makes these
equivalencies explicit.\\[5pt]

\noindent
Newly available astrophysics courses have been added as electives: PHY
156 (Introduction to Cosmology) and PHY 158 (Formation of Galaxies \& Cosmic Structures).  The
requirement for one additional upper-division physics course has
been replaced by a second elective, so that every upper division
course in the major is chosen from a coherent list.  PHY 117 and 118
are now available as electives.\\[5pt]

\noindent
{\bf Prompt c.} {\it Describe the timing of and implementation plan for the revisions. Please be sure to address all impacts to students and catalog rights, including:}\\
{\bf i.} {\it How many existing students in the major will be affected by the proposed change?}\\
{\bf ii.} {\it Are any mitigation measures needed to prevent delays in time to degree for students already enrolled in the major or minor?}\\

\noindent
We have about 2 physics AB majors per year, and typically they are double majors with another
discipline.  Mostly these changes add additional options for students.  PHY 40 is already being offered and was previously a hidden prerequisite of PHY 104A; it is
now required explicitly.  The additional upper division physics course that students already take
must now come from the list of electives.\\[5pt]

\noindent
There is a clear equivalency of requirements that will allow us to
handle special cases as needed:
\begin{center}
\begin{tabular}{ll}
Old Requirement & Satisfied By \\
\hline
PHY 9HE & Not required \\
PHY 102 & PHY 40 \\
\end{tabular}
\end{center}

\noindent
{\bf Prompt d.} {\it Each proposal should contain revisions for only
  one major or minor program. Occasionally, multiple related proposals
  need to be approved in sequence to avoid disruption to
  enrollments. Please indicate this in cover letters, but submit
  separate proposals.}\\

\noindent
This proposal is focused exclusively on the AB Physics major.  We expect
to submit updates to the BS Physics and BS Applied Physics majors
separately, but these can be approved independently.

\section*{3. Consultation with Affected Units (e.g., Departments, Programs, Centers)}

{\bf Prompt a.} {\it Describe overlap or conflict with existing UC Davis majors or minors.}\\

\noindent
The closest major to an AB Physics offered by another department is the
BS in Applied Mathematics.  Students in that major might select PHY 9A
and 9B plus one enrichment course that could include an upper-level
physics requirement such as 105A or 104A.  The focus is much more
heavily on mathematics, and overall this is not a significant overlap.
In terms of upper division course work, this is four units in physics
compared to 36 units (minimum) for our AB majors.\\[5pt]

\noindent
{\bf Prompt b.} {\it You must consult with affected units and include letters from such units stating their feedback. Proposals may still be submitted even if affected units are not supportive. The Davis Division will only review proposals that include:}\\
{\bf i.} {\it Feedback from affected units; or}\\
{\bf ii.} {\it If affected units choose not to provide feedback, documentation of reasonable attempts to receive this feedback.}\\

\noindent
We have included a letter of support from the Department of Mathematics.  We are adding options for
additional new variants of required math courses, as requested by the math department.

\section*{4. Curricular Structure}

\noindent
{\bf Prompt a.} {\it Revisions to the curricular structure should be submitted in track changes (deletion and addition of courses and changes to total required units should be clear).}\\

\noindent
We have submitted the updated curricular structure in track changes.\\[5pt]

\noindent
{\bf Prompt b.}{\it Provide letters of support from units whose courses are added to the program in this revision.}\\
{\it i. From Chairs/Directors: Confirm courses have capacity to accommodate frequency of offerings and number of projected students; comment on any other levels of unit/program support required.}\\[5pt]

\noindent
We have included a letter of support from the Department of Mathematics.\\[5pt]

\noindent
{\bf Prompt c.}  {\it Note that as per COCI policy, UGC acknowledges all in-person, virtual (V), and
  hybrid (Y) versions of a course number to be interchangeable with respect to major and
  prerequisite requirements. In such cases where differentiation is desired (one version accepted
  toward the degree and another not), the program should contact COCI about required adjustments to
  course numbering.}\\

\noindent
We have no need to make such a request at this time.\\[5pt]

\noindent
{\bf Prompt d.}
{\it Consider the impact of prerequisites to required courses that are not explicitly listed among major requirements (“hidden prerequisites”) on time to degree.}\\

\noindent
We have been careful to remove any hidden prerequisites.  We have added explicit requirements for
prerequisites as needed for all required courses.\\[5pt]

\noindent
Under the current requirements, PHY 40 is a hidden prerequisite (via
PHY 104A).  The proposed requirements make it explicitly required;
see the Required Units section below.

\section*{Required Units}

\noindent
In the UC Davis course catalog, the minimum units required for the AB
Physics major is listed as 81 units.  However, PHY 104A, a required
course, has PHY 40 as a prerequisite, which the current requirements
do not list.  Taking PHY 40 adds three units, less one unit for PHY
102, which is waived when PHY 40 is taken.  {\bf The actual minimum
  number of units required to complete the current AB Physics
  requirements is 83 units, despite what the catalog claims.}  Now that
MAT 22A and 22B are four-unit courses, this minimum is 85 units, five
units above the cap.\\[5pt]

\noindent
This proposal changes this minimum to 86 units.  Here's an
accounting of which changes impact the minimum:
\begin{itemize}
\item Drop the additional fixed-unit upper-division physics course (-3)
\item Require a second elective (+4)
\end{itemize}
Making PHY 40 explicitly required, and adding PHY 117, 118, 156, and
158 to the electives, do not change the minimum.  The proposal
requires 36 upper division units, the minimum of L\&S Regulation
34(A)(1).\\[5pt]

\noindent
This proposal therefore requires an exemption under L\&S Regulation
33(D).  For consistency across our majors, we state the exemption as a
general rule:
\begin{quote}
  Lower division courses offered outside the major department (Physics and
  Astronomy), and required for the major (including as prerequisites),
  count at 75\% of their units toward the limit of 80 units
  (Regulation 33(B)), reducing the count by at most 8 units.  The major
  may not exceed 88 units.
\end{quote}
For the AB, the 24 units of mathematics count as 18, and the major
counts as 80 units.  The exemption therefore allows the AB to require
86 units, which Regulation 33(D) permits.\\[5pt]

\noindent
{\bf Lower division prerequisites have grown, and the limit is no
  longer workable.}  MAT 22A and 22B are now four-unit courses, and
PHY 40 is now a prerequisite of PHY 104A.  No path through the AB both
reaches the 36 upper division units required by Regulation 34(A)(1)
and stays within 80 units:  the preparatory coursework alone requires
50 units.  The current AB already requires 85 units, and only 35 upper
division units.\\[5pt]

\noindent
{\bf The limit exists to preserve breadth, and the proration applies
  only to coursework that provides it.}  Regulation 33(D) anticipates
exceptions ``particularly when the academic major includes
considerable breadth.''  Coursework in other departments is breadth,
and only that coursework is prorated:  every course in the major
department, and all upper division coursework, counts in full.  The
reduction is limited to 8 units, so that the exemption cannot be used
to grow the major without bound.\\[5pt]

\noindent
The exemption is presented in full in Section 2 of the longer write-up.\\[5pt]

\noindent
{\bf The physics department stands behind the minimum unit
  calculations presented here, and in our full proposal, as we have
  had independent confirmation of our tallies several times now.}
These unit minimums are surprisingly challenging to calculate
correctly.  A more detailed tally of required units under both the
current and proposed requirements is available in Appendix C of the
longer write-up.

\end{document}
